\section{问题三：多人拥挤博弈与在线反馈}

\subsection{多人规则的精确恢复}

原题中的多人系数以旧版 Word 的 WMF 公式保存，普通文本抽取会遗漏。通过原始页面渲染核实：同一天同一有向边上有 $k\ge2$ 人时，每人移动消耗系数由2变为 $2k$；同矿 $k$ 人的消耗系数仍为3，每人收入为 $B/k$；同村至少两人购买时，单价由基准价2倍升为4倍。第五关 $n=2$，第六关 $n=3$。

对玩家 $i$，给定联合行动 $\bm a_t$，移动消耗倍数为
\begin{equation}
\lambda^i_t=
\begin{cases}
2k_{t,e},&a^i_t=\text{沿有向边 }e\text{ 移动},\ k_{t,e}\ge2,\\
2,&a^i_t=\text{移动且 }k_{t,e}=1,\\
3,&a^i_t=\text{挖矿},\\
1,&a^i_t=\text{停留}.
\end{cases}
\end{equation}
记$k_{t,m}$为同矿挖矿人数，$k_{t,v}^{B}$为同村同日购买人数。玩家$i$的资源和现金转移可统一写成
\begin{align}
r_t^i&=r_{t-1}^i+b_t^{r,i,-}+b_t^{r,i,+}
-\lambda_t^ic_{z_t}^r,\qquad r\in\{w,f\},
\label{eq:multi-resource}\\
m_t^i&=m_{t-1}^i+\frac{B}{k_{t,m}}I_{t,i}^{M}
-\rho(k_{t,v}^{B})\sum_{r}p_rb_t^{r,i},
\label{eq:multi-cash}\\
\rho(k)&=
\begin{cases}4,&k\ge2,\\2,&k=1.\end{cases}
\label{eq:village-congestion}
\end{align}
其中无挖矿时令$I_{t,i}^{M}=0$。联合状态$\bm S_t=(S_t^1,\ldots,S_t^n)$不再能按玩家独立转移，因为$\lambda_t^i$、矿山收益和采购价格都由联合行动决定。令$\bm\Phi$为“先分组、再结算”的联合映射，则
\begin{equation}
\bm S_t=\bm\Phi(\bm S_{t-1},z_t,\bm a_t,\bm b_t),
\label{eq:joint-transition}
\end{equation}
这也是多人模拟器逐日执行的数学定义。

这使玩家收益 $J_i(\pi_i,\pi_{-i})$ 显式依赖其他玩家策略。Nash均衡要求
\begin{equation}
J_i(\pi_i^*,\pi_{-i}^*)\ge J_i(\pi_i,\pi_{-i}^*),\quad
\forall i,\ \forall\pi_i,
\label{eq:nash}
\end{equation}
与非合作博弈的经典定义一致\cite{nash1951}。

\subsection{第五关：完整最佳响应动态规划}

天气和对手计划固定时，自身最优响应仍可用动态规划求解。状态为
\[
(t,l_t,W_t,F_t),
\]
其中$W_t,F_t$是截至第$t$日累计需要购买的资源。给定对手第$t$日的位置—行动$(l_t^{-i},a_t^{-i})$，自身行动$a$产生确定增量
\begin{equation}
(\Delta W_t,\Delta F_t,\Delta R_t)
=g_t(l,a;l_t^{-i},a_t^{-i},z_t),
\label{eq:best-response-increment}
\end{equation}
其中同向边把移动倍数由2改为4，同矿把收入由200元改为100元。定义$D_t(l,W,F)$为到达该标签时可得的最大累计矿山收入，则标签递推为
\begin{equation}
D_t(l',W',F')=max_{a\in\mathcal A_t(l)}
\left\{D_{t-1}(l,W,F)+\Delta R_t(a)\right\},
\label{eq:best-response-dp}
\end{equation}
其中$l'=L(l,a)$、$W'=W+\Delta W_t$、$F'=F+\Delta F_t$。只保留满足
\begin{equation}
3W'+2F'\le C,\qquad p_wW'+p_fF'\le M_0,\qquad
d(l',v_T)\le T-t
\label{eq:best-response-pruning}
\end{equation}
的标签。相同$(l',W',F')$只保留$D_t$最大的路径，因为它对未来可行域无影响而终值更高。终点标签的价值为
\begin{equation}
J_i=M_0-p_wW_\tau-p_fF_\tau+D_\tau(v_T,W_\tau,F_\tau).
\label{eq:best-response-terminal}
\end{equation}
由有限日期、有限节点和整数容量可知标签集有限，故\cref{eq:best-response-dp}枚举了固定对手下的完整策略空间，而不是局部换路启发式。

从单人最优计划出发交替计算最佳响应，在30轮上限内收敛；随后固定对手，对双方各再做一次完整最佳响应。定义偏离增益
\begin{equation}
\Delta_i=\max_{\pi_i}J_i(\pi_i,\pi_{-i}^*)-J_i(\pi_i^*,\pi_{-i}^*).
\label{eq:deviation-gain}
\end{equation}
只有当$\Delta_1=\Delta_2=0$时才接受纯策略Nash均衡。结果见\cref{tab:level5-eq}。

\begin{table}[H]
\centering
\caption{第五关非对称纯策略Nash均衡}
\label{tab:level5-eq}
\begin{tabularx}{0.94\textwidth}{C{1.2cm} C{1.8cm} Y C{1.2cm} C{1.8cm} C{2.0cm}}
\toprule
玩家 & 初始水/食物 & 行动序列 & 到达日 & 终值/元 & 最大偏离增益/元\\
\midrule
1 & 27/33 & $1\to4$，停留，$4\to6\to13$ & 4 & \textbf{9535} & 0\\
2 & 30/34 & $1\to5\to6\to13$ & 3 & \textbf{9510} & 0\\
\bottomrule
\end{tabularx}
\end{table}

第2日是高温，玩家1在4区停留，把两个移动日推迟到第3、4日晴朗；玩家2从5区走另一条首段路线。两人的同日有向边没有重叠。低基础收益使绕矿不具吸引力，因此均衡核心是“天气错峰+空间分流”，而不是抢矿。

\begin{figure}[H]
  \centering
  \includegraphics[width=0.98\textwidth]{level5_equilibrium.pdf}
  \caption{第五关均衡的空间路径、行动时序与偏离复核}
  \label{fig:level5}
\end{figure}

\subsection{第六关：三人联合反馈模拟}

第六关每天先揭示天气，再由三名玩家根据上一日公开的联合状态选择行动。联合反馈策略是映射
\begin{equation}
\bm\pi_t:\ (\bm S_0,z_1,\bm a_1,\ldots,\bm S_{t-1},z_t)
\longmapsto \bm a_t,
\label{eq:joint-policy}
\end{equation}
其后由\cref{eq:joint-transition}一次性结算。模拟器严格按行动前村庄购买分组、同向边分组、同矿分组、资源扣除、收益分摊、行动后购买、失败和终点结算的顺序执行。为分离不同外部性，比较四种策略：

\begin{enumerate}
  \item \textbf{镜像挖矿：}三人复制单人矿山阈值策略；
  \item \textbf{分路直达：}两人分别走上侧和左侧8日最短路，第三人错峰1个非沙暴日后跟随；
  \item \textbf{拥挤感知非合作：}先分路到矿山，在场且资源安全者都加入挖矿；
  \item \textbf{合作轮换：}先分路到矿山，每天最多安排一人挖矿，作为可协调对照。
\end{enumerate}

先在IID假设下建立效率基准。分路直达的初始水食从180到220箱、步长5搜索，共9个候选。设装载参数$q$表示每人初始水、食物均为$q$箱，训练选择为
\begin{equation}
q^*=\arg\max_{q\in\{180,185,\ldots,220\}}
\widehat\mu_{\mathrm{adv}}(q)
\quad\text{s.t.}\quad
\widehat p_{\mathrm{base}}(q)\le0.002,
\ \widehat p_{\mathrm{adv}}(q)\le0.005.
\label{eq:level6-load-selection}
\end{equation}
训练阶段基准和不利分布各3000局；入选$q^*=185$。测试阶段换用种子20260921，每种分布20000局，即60000个玩家观测。因为同一局三人的失败相关，玩家级区间用于描述个体风险，同时另行计算局级区间，避免把60000个观测误当作完全独立的局。

\begin{figure}[H]
  \centering
  \includegraphics[width=0.96\textwidth]{level6_tradeoff.pdf}
  \caption{第六关四类反馈策略跨天气分布的失败风险与总体终值}
  \label{fig:level6-tradeoff}
\end{figure}

\begin{table}[H]
\centering
\caption{第六关分路直达策略的独立样本外结果}
\label{tab:level6-direct}
\begin{tabular}{lrrrr}
\toprule
天气模型 & 玩家失败率 & 95\% Wilson区间 & 总体均值/元 & 5\%分位/元\\
\midrule
低沙暴 & 0 & [0, 0.0064\%] & 7820.3 & 7612.5\\
基准 & 0.0167\% & [0.0091\%, 0.0307\%] & \textbf{7767.6} & 7542.5\\
不利 & 0.415\% & [0.3666\%, 0.4697\%] & 7637.4 & 7395.0\\
\bottomrule
\end{tabular}
\end{table}

基准分布下，镜像挖矿、拥挤感知非合作、合作轮换的玩家失败率分别为97.865\%、12.362\%和6.978\%；总体均值分别为142.6、7122.8和8326.6元。合作轮换均值高于直达，却远超0.2\%风险预算；故IID层面的风险受控基准是分路直达，下一节再检验其面对天气持续性的稳定性。

\subsection{角色条件装载与分布鲁棒策略}

185/185箱是IID模型下的效率点，但连续坏天气会使它越过风险预算。沿用\cref{eq:weather-ambiguity-set}的模糊集，将最短路角色按出发前自愿错峰天数$d$分类。对$d\in\{0,1,2\}$，在水、食物180--240箱、步长5的二维网格$\mathcal L$上校准
\begin{equation}
(w_d^*,f_d^*)=
\arg\max_{(w,f)\in\mathcal L}
\min_{\mathbb P\in\mathcal P}\widehat{\mathbb E}_{\mathbb P}[J(w,f,d)]
\quad\text{s.t.}\quad
U_W\!\left(X_{\mathbb P};N\right)\le\alpha(p),
\ \forall\mathbb P\in\mathcal P,
\label{eq:role-conditional-load}
\end{equation}
其中$U_W$是\cref{eq:wilson-center,eq:wilson-halfwidth}给出的失败率Wilson上界，而非失败率点估计。没有边拥挤时，八个移动日加$d$个非沙暴停留日的资源需求可直接写成
\begin{align}
R_w(z_{1:\tau};d)&=\sum_{t=1}^{\tau}\lambda_t(d)c^w_{z_t},\qquad
R_f(z_{1:\tau};d)=\sum_{t=1}^{\tau}\lambda_t(d)c^f_{z_t},
\label{eq:direct-role-requirement}\\
\lambda_t(d)&=
\begin{cases}
1,&z_t=\text{沙暴或当日用于自愿错峰},\\
2,&\text{当日沿最短路移动},
\end{cases}
\end{align}
到达所需的第八个非沙暴移动日为$\tau$。若$R_w>w$或$R_f>f$则失败记0；否则终值为
\begin{equation}
J=10000-\frac12p_ww-\frac12p_ff
-\frac12p_wR_w-\frac12p_fR_f,
\label{eq:direct-role-value}
\end{equation}
因为未消耗库存按半价回收。该式也解释了为什么在通过风险约束后，继续加装会单调降低期望终值。

每个模糊集点使用30000条校准序列。$d=0$的最优可行装载为215/215箱，$d=1$为225/230箱；$d=2$在1200千克容量内没有同时通过两档风险预算的装载。因此最终剖面为两名玩家立即从不同方向出发、各装215/215箱，第三名玩家装225/230箱并错峰一个非沙暴日。公开随机置换三个角色后，每个人承担相同的事前混合风险，但角色条件装载还额外保证每个固定角色单独过线。

\begin{table}[H]
\centering
\caption{第六关分布鲁棒剖面的独立样本外结果}
\label{tab:level6-robust}
\begin{tabular}{lrrrr}
\toprule
边际天气 & $\rho$ & 玩家失败率 & 固定角色最大Wilson上界 & 总体均值/元\\
\midrule
基准 & 0 & 0 & 0.0128\% & 7509.6\\
基准 & 0.35 & 0.0600\% & 0.1030\% & 7499.4\\
不利 & 0 & 0.0056\% & 0.0390\% & 7410.0\\
不利 & 0.35 & 0.3400\% & 0.4344\% & \textbf{7383.3}\\
\bottomrule
\end{tabular}
\end{table}

表中每档30000局。“玩家失败率”按三角色合并统计；验收列则分别对三个固定角色各用30000个独立天气样本计算Wilson上界，再取最大值，避免把同局三人的相关结果误当作90000个独立Bernoulli观测。模糊集内最坏固定角色上界0.4344\%仍低于0.5\%预算。

\begin{figure}[H]
  \centering
  \includegraphics[width=0.98\textwidth]{robust_load_calibration.pdf}
  \caption{第六关错峰角色的二维装载可行域与分布鲁棒入选点}
  \label{fig:robust-load-calibration}
\end{figure}

\subsection{受约束最佳响应预言机与可利用度}

5乘5网格从1到25的最短路需要4次向右和4次向下，因此共有
\begin{equation}
|\mathcal R|=\binom{8}{4}=70
\label{eq:shortest-route-count}
\end{equation}
条。构造偏离类$\Pi_i^{\mathcal C}$：玩家$i$可在70条最短路中任意换路，并选择$d=0$或1；装载使用\cref{eq:role-conditional-load}为该$d$校准的风险最优装载。$d=2$因不存在风险可行装载而剔除。于是每名玩家有$70\times2=140$个可审计偏离，而对手保持推荐剖面。

令分布鲁棒收益$\underline J_i(\pi)=\min_{\mathbb P\in\mathcal P}\mathbb E_{\mathbb P}[J_i(\pi)]$，定义受约束可利用度
\begin{equation}
\epsilon_i^{\mathcal C}=
\max_{\pi_i'\in\Pi_i^{\mathcal C}:\,U_W\le\alpha}
\underline J_i(\pi_i',\pi_{-i}^*)
-\underline J_i(\pi_i^*,\pi_{-i}^*).
\label{eq:class-exploitability}
\end{equation}
预言机先在经验最不利角点“不利边际、$\rho=0.35$”用1500条配对序列完整筛选，再在$\mathcal P$每个点各用30000条新序列复核入选偏离。三个角色的最佳响应都只是换到另一条不发生同向边重合的等价最短路，样本外收益与原角色逐样本相同，故$\epsilon_i^{\mathcal C}=0$。这把原先的“比较四种策略”推进到显式偏离检验；但$\Pi_i^{\mathcal C}$不含绕矿、村庄补给及任意历史依赖策略，因此结论是类约束$\epsilon$-均衡，而非整个随机动态博弈的Markov完美均衡。

\begin{figure}[H]
  \centering
  \includegraphics[width=0.90\textwidth]{level6_exploitability.pdf}
  \caption{第六关推荐剖面与受约束最佳响应的样本外收益比较}
  \label{fig:level6-exploitability}
\end{figure}

\subsection{外部性分解与角色公平}

\cref{fig:level6-mechanism}给出基准天气的机制量。镜像挖矿每局产生19.44个共享边玩家日，三人往往尚未到矿山就因6倍移动消耗失败。拥挤感知非合作把共享边降至4.02，却产生8.74个共享矿山玩家日；合作轮换消除共享矿山，但等待时间抬高尾部失败。分路直达两类共享玩家日均为0。

\begin{figure}[H]
  \centering
  \includegraphics[width=0.98\textwidth]{level6_mechanism.pdf}
  \caption{第六关多人策略的共享边、共享矿山与失败率分解}
  \label{fig:level6-mechanism}
\end{figure}

分布鲁棒装载下，第三名玩家因错峰与更高安全库存，在基准IID分布平均比前两名少155.2元；不利边际、$\rho=0.35$时角色差为165.9元。设三个路线角色的期望终值为$\mu_1,\mu_2,\mu_3$，第0天生成与身份无关的均匀随机置换$\sigma$，把玩家$i$分配到角色$\sigma(i)$。则对任意玩家$i$，
\begin{equation}
\mathbb E[J_i]
=\sum_{r=1}^3\Pr\{\sigma(i)=r\}\mu_r
=\frac13(\mu_1+\mu_2+\mu_3)=\overline\mu.
\label{eq:exante-fairness}
\end{equation}
所以随机置换不改变三条路线的联合占用和系统总收益，却使每名玩家获得相同的事前期望。固定角色的事后收益仍可能不同，这里主张的是事前公平而非逐局完全相等。

\begin{resultbox}{问题三结论}
第五关的9535/9510元组合通过完整最佳响应检验，属于精确纯策略均衡\evidence{E3}。第六关在IID下以185/185箱为效率基准；最终分布鲁棒推荐采用215/215、215/215与225/230箱的异路错峰剖面，并在每名玩家140个路线—错峰偏离类中得到零可利用度\evidence{E4}。该结论不外推为完整Markov完美均衡。
\end{resultbox}
